\documentclass[../main.tex]{subfiles}
\begin{document}
\paragraph{Hebbian learning rule states that neurons that fire together wire together. This is a fundamental rule in used for training of neural networks but on it's own it is only useful for a single layered neural network.}
\subsection{Assumptions}
\paragraph{We assume the following when implementing the hebbian learning rule.}
\begin{itemize}
  \item The netowork only has one input layer of size $M$ and one output layer of size $N$.
  \item There is a weight matrix of size $[M \times N]$.
  \item There is only one bias vector of size $N$.
  \item The activation function is applied only on the output layer.
\end{itemize}
\subsection{Equations}
\subsubsection{Forward}
\text{Let's say we have an input vector of size $M$ given as:}
$$ X = \begin{bmatrix}
  x_1 & x_2 \cdots x_M
\end{bmatrix}$$
\text{and we have weight matrix of size $M \times N$ given as:}
  $$ W = \begin{bmatrix}
  w_{11} & w_{12} \cdots w_{1N} \\
  w_{21} & w_22 \cdots w_{2N} \\
   \vdots & \\
  w_{M1} & w_{M2} \cdots w_{MN}\\
\end{bmatrix}$$
\text{and we have bias in a vector size $N$ given as:}
$$ b = \begin{bmatrix}
  b_1
  \cdots
  b_N
\end{bmatrix}$$
\text{using these we can write the forward pass as:}
$$ y = activation(X \cdot W + b)$$
\text{here $y$ is the predicted output of the neural network.} \\
\text{$y$ is vector of size $N$}
\text{"." is the matrix multiplication.}
\subsubsection{Updates}
\paragraph{Hebbian learning rule has only two updates, one is for weights and other is for biases. These grow infinitely so training for longer intervals can make model useless.}
\begin{align}
  \Delta w &= \eta y.x \\
  \Delta b &= \eta y
\end{align}
The $\eta$ is the learing rate. \\
y is the output vector. \\
x is input vector.
\subsubsection{Insights}
\begin{itemize}
  \item Hebbian learning rule is unsupervised.
  \item The dataset should be divided into positive and negative category becuase it will help in driving the weights up or down.
\end{itemize}
\end{document}
